% !PRINTED 587
\weberpage{587}
Andererseits ist $f_\nu(\tau)\,\psi_\mu(\sigma)$ ein ganzes Functional des Körpers
$\Omega$, und folglich nach dem vorhin schon angewandten Satze \S.\,158, 4.:

\weq{16}{
\begin{aligned}
S_\Omega \left( \frac{f_\nu(\tau)\,\psi_\mu(\sigma)}{F'(\tau)} \right) &= b_{\mu,\nu} &
\mu &= 0, 1, \ldots, m-1\\
& & \nu &= 0, 1, \ldots, n-1
\end{aligned}
}
ein ganzes rationales Functional. Multiplicirt man diese Gleichung mit
$\tau^{n-\nu-1}\sigma^{m-\mu-1}$ und nimmt die Summe über alle $\mu$ und $\nu$, so folgt,
wie vorhin, da für zwei von $\sigma, \tau$ verschiedene Wurzeln $\sigma_k, \tau_i$ von
$\psi=0, f=0$
\[
\begin{aligned}
\sigma^{m-1}\psi_0(\sigma_k) + \sigma^{m-2}\psi_1(\sigma_k) + \cdots + \psi_{m-1}(\sigma_k) &= 0,\\
\tau^{n-1} f_0(\tau_i) + \tau^{n-2} f_1(\tau_i) + \cdots + f_{n-1}(\tau_i) &= 0,
\end{aligned}
\]
und
\[
\begin{aligned}
\sigma^{m-1}\psi_0(\sigma) + \sigma^{m-2}\psi_1(\sigma) + \cdots + \psi_{m-1}(\sigma) &= \psi'(\sigma),\\
\tau^{n-1} f_0(\tau) + \tau^{n-2} f_1(\tau) + \cdots + f_{n-1}(\tau) &= f'(\tau)
\end{aligned}
\]
ist, die Gleichung
\weq{17}{
\frac{f'(\tau)\,\psi'(\sigma)}{F'(\tau)} = \sum^{\mu,\nu} b_{\mu,\nu}\, \tau^{n-\nu-1}\sigma^{m-\mu-1}.
}

Also ist auch $f'(\tau)\,\psi'(\sigma)$ durch $F'(\tau)$ theilbar und unser Theorem (8)
daher bewiesen.

Bezeichnen wir mit $G_\Omega, G_R$ die Grundideale der Körper $\Omega, R$, mit
$\mathfrak{G}_R$ das relative Grundideal von $\Omega$ in Bezug auf $R$, so kann man dem
bewiesenen Theorem auch den Ausdruck geben:
\weq{18}{
G_\Omega = \mathfrak{G}_R\, G_R.
}
Es ist klar, wie sich dieses Theorem verallgemeinern lässt: Es sei
\[
\Omega_1, \Omega_2, \Omega_3, \ldots
\]
eine Reihe von Körpern, deren jeder die folgenden als Theiler enthält. Es sei ferner
$\mathfrak{G}_{i,k}$ das Grundideal von $\Omega_i$ in Bezug auf $\Omega_k$ ($k > i$),
dann ist
\weq{19}{
\mathfrak{G}_{i,k} = \mathfrak{G}_{i,i+1}\,\mathfrak{G}_{i+1,i+2} \cdots \mathfrak{G}_{k-1,k}.
}
Die tiefere Bedeutung dieser Grundideale wird sich in einem der nächsten Paragraphen noch
deutlicher herausstellen.

Wir wollen aber schon hier auf eine merkwürdige Erscheinung aufmerksam machen, dass es
nämlich vorkommen kann, dass $f'(\tau)$ eine Einheit ist, was nach dem
\gesperrt{Minkowski'schen} Satze nicht möglich ist, wenn $R$ der Körper der rationalen
Zahlen ist. In diesen Fällen, von denen die Theorie der com-
%CONT>>
